\documentclass[10pt]{article}
\usepackage[margin=0.62in]{geometry}
\usepackage[T1]{fontenc}
\usepackage{lmodern}
\usepackage{microtype}
\usepackage{hyperref}
\usepackage{enumitem}
\usepackage{booktabs}
\usepackage{array}
\usepackage{xcolor}
\usepackage{titlesec}
\usepackage{amsmath}

\hypersetup{colorlinks=true,urlcolor=blue!55!black,linkcolor=blue!55!black}
\setlength{\parindent}{0pt}
\setlength{\parskip}{3pt}
\setlist[itemize]{leftmargin=*,nosep,topsep=2pt}
\setlist[enumerate]{leftmargin=*,nosep,topsep=2pt}
\titlespacing*{\section}{0pt}{7pt}{3pt}
\titlespacing*{\subsection}{0pt}{5pt}{2pt}

\title{\vspace{-1.2em}\textbf{Docling Formula Adapter: Premise, Field Map, and Gap}\\[-0.2em]
\large A two-page conversation brief for Taru}
\author{}
\date{\vspace{-1.6em}June 13, 2026}

\begin{document}
\maketitle
\vspace{-1.5em}

\section{Premise}

We like \textbf{Docling} because it is already shaped like infrastructure, not a one-off model demo: it parses documents into a structured representation, supports layout/table/formula/code enrichment, exposes Python and CLI surfaces, and fits downstream RAG/data pipelines. The interesting opportunity is not ``OCR but with another VLM.'' It is \textbf{verifiable structured extraction from scientific documents}, beginning with formulas.

The seed idea: build or adapt a small Docling-compatible formula recognizer that converts formula crops/pages into \LaTeX, preserves location/provenance, exposes confidence and validation metadata, and later supports a formal subset that can be checked in Lean.

\textbf{Important framing:} Lean does not validate arbitrary \LaTeX. \LaTeX is display syntax; Lean is typed formal math. The practical stack is:
\[
\text{image} \rightarrow \text{LaTeX} \rightarrow \text{canonical LaTeX} \rightarrow \text{render/visual checks} \rightarrow \text{optional typed Lean check}
\]

\section{What The Field Has Done Well Since 2023}

\textbf{1. Scientific PDF-to-markup became real.} Nougat showed that academic PDFs can be converted into markup with formulas, attacking the ``PDF loses semantics'' problem directly. It proved demand and set the scientific-document framing, but it is academic-document-specialized and not a full general document system.

\textbf{2. General document parsers got much stronger.} Docling, MinerU, Dolphin, PaddleOCR-VL, and olmOCR pushed the field toward full document conversion: reading order, layout, tables, formulas, charts, lists, scanned pages, and Markdown/HTML-like outputs. The best systems now combine learned models with careful preprocessing/postprocessing and benchmark discipline.

\textbf{3. Small/compact VLMs became credible.} SmolDocling is the key Docling-native example: a 256M VLM generating DocTags with layout and content. GOT-OCR2.0 is another strong signal: a 580M unified OCR-2.0 model for text, formulas, tables, charts, geometry, and formatted outputs. PaddleOCR-VL targets a compact 0.9B multilingual parser. The lesson is that good visual tokenization and data curation can beat brute force for document tasks.

\textbf{4. Synthetic data works when it is aligned and checkable.} SynthFormulaNet and SynthCodeNet are large, controlled image-to-text corpora. ChartNet goes further by aligning chart image, plotting code, source table, summary, and QA. olmOCR 2 is especially relevant because it uses unit-test-like verifiable rewards for document OCR, with gains in formulas, tables, and multi-column layouts.

\textbf{5. The frontier is moving from OCR to structured visual code.} GOT, DeepSeek-OCR, MOCR/dots.mocr, and chart/table systems all point the same way: document elements are not just pixels to transcribe; they are visual programs to recover. Formulas are one of the purest cases of this problem.

\section{What Is Already Done Super Well}

\begin{tabular}{p{0.24\linewidth}p{0.68\linewidth}}
\toprule
\textbf{Area} & \textbf{Strong existing work} \\
\midrule
Document pipeline & Docling gives a clean, extensible document conversion stack. We should build with it, not around it. \\
Formula/code baseline & CodeFormulaV2 already handles image crops of formulas/code and is trained on SynthFormulaNet/SynthCodeNet. \\
Full-page conversion & SmolDocling, olmOCR, MinerU, Dolphin, and PaddleOCR-VL show strong page-level extraction and reading-order work. \\
Efficient VLM design & Vary-toy, SmolDocling, GOT, DeepSeek-OCR, and PaddleOCR-VL show compact models can work if visual tokens/data are good. \\
Synthetic supervision & SynthFormulaNet, SynthCodeNet, ChartNet, and olmOCR-style synthetic unit tests show scalable supervision is possible. \\
\bottomrule
\end{tabular}

\section{What Still Does Not Feel Solved}

\textbf{Formula quality is not just exact-match OCR.} Two \LaTeX strings may render identically; one string may compile but mean the wrong thing; another may be visually right but semantically under-specified. We need layered evaluation: canonical string, compile/render, visual similarity, and optional semantic checks.

\textbf{Real scientific PDFs are messy.} Synthetic formulas are clean. Real papers have equation numbers, line breaks, multi-line alignments, low-resolution scans, font issues, cropped symbols, ambiguous glyphs, theorem context, and notation defined in prose.

\textbf{Model outputs rarely carry trust metadata.} A downstream agent needs to know: did this formula compile, did it visually match the crop, was it low confidence, did alternate candidates exist, and is this formula eligible for Lean validation?

\textbf{Formal validation is underused and often misunderstood.} Lean can be powerful, but only after a typed semantic parse. The right move is a narrow high-value Lean subset, not pretending every formula crop can be auto-formalized.

\textbf{There is no obvious ``formula adapter benchmark'' for Docling users.} Existing systems report broad document metrics. A hard, curated, formula-centered benchmark with validation traces would be useful even before model training.

\section{Proposed Contribution}

The gap is not ``make another OCR model.'' A potentially useful contribution is a \textbf{review/benchmark}: compare VLM and non-VLM formula extraction approaches for Docling-like systems, especially on-device, and identify where CodeFormulaV2 or current SOTA is best versus where another approach may be cheaper, faster, or more reliable.

\begin{quote}
\textbf{Target artifact:} a Docling-native formula extraction benchmark and adapter study with validation traces: image-to-\LaTeX, canonicalization, compile/render checks, visual equivalence, time-to-parse, memory footprint, layout failure taxonomy, and a small Lean-checkable subset.
\end{quote}

\section{Targets We Should Optimize}

\begin{tabular}{p{0.24\linewidth}p{0.68\linewidth}}
\toprule
\textbf{Target} & \textbf{Metric audiences may care about} \\
\midrule
On-device viability & Model size, peak RAM, CPU/MPS latency, cold start, no cloud dependency. \\
Faster parsing & Total PDF parse time, seconds/page, formulas/second, overhead versus Docling without formula enrichment. \\
Higher accuracy & Normalized \LaTeX exact match, symbol error rate, structure/tree similarity, render similarity, compile rate. \\
Layout handling & Formula detection recall, bbox quality, inline/display split, multi-line equations, equation numbers, reading order. \\
Trust metadata & Confidence calibration, abstention quality, alternate candidates, compile/render pass, visual-match score, Lean-valid subset. \\
\bottomrule
\end{tabular}

The practical bar is a Pareto curve, not one score: \textbf{accuracy vs latency vs memory vs trust}. Taru's point is central: an accurate system that takes 30 minutes for 10 pages is not useful for real Docling workflows.

\section{Possible First Milestone}

\begin{enumerate}
\item Sample and inspect SynthFormulaNet plus a small set of real math-heavy PDFs.
\item Build a validator: parse target, render \LaTeX, compare to source crop, tag failures.
\item Run CodeFormulaV2, SmolDocling/Granite-Docling, one larger VLM baseline, and one non-VLM/specialized formula OCR baseline if available.
\item Create a hard 500--2,000 example benchmark stratified by formula type.
\item Report accuracy, parse time, memory, and layout failures on 10-page, 50-page, and math-heavy stress PDFs.
\item Wrap best baseline/cascade as a Docling enrichment adapter with validation metadata.
\item Add a tiny Lean-valid pilot: elementary algebra/calculus/theorem statements with explicit assumptions.
\end{enumerate}

\section{Questions For Taru}

\textbf{Decision questions:} benchmark vs adapter vs validation tool? arXiv/math PDFs vs textbooks vs scanned papers? local speed vs top accuracy vs trust metadata? which Lean slice is realistic first? Paper shortlist for follow-up: CodeFormulaV2, SynthFormulaNet, SmolDocling, Nougat, GOT-OCR2.0, DeepSeek-OCR/2, olmOCR 2, DocTron-Formula, ChartNet.

\end{document}
